\documentclass[a4paper,11pt]{article}

\usepackage[a4paper,margin=1.8cm]{geometry}
\usepackage{fontspec}
\usepackage{polyglossia}
\setmainlanguage{russian}
\setmainfont{Arial}
\setmonofont{Arial}

\usepackage{amsmath}
\usepackage{array}
\usepackage{booktabs}
\usepackage{colortbl}
\usepackage{xcolor}
\usepackage{tabularx}
\usepackage{longtable}
\usepackage{enumitem}
\usepackage{fancyhdr}
\usepackage{hyperref}
\usepackage{microtype}

\definecolor{accent}{HTML}{2457A6}
\definecolor{lightblue}{HTML}{EAF1FB}
\definecolor{darkgray}{HTML}{39424E}

\hypersetup{
  colorlinks=true,
  linkcolor=accent,
  urlcolor=accent,
  citecolor=accent,
  pdftitle={LLM Classification Fine-tuning Technical Report}
}

\renewcommand{\arraystretch}{1.32}
\setlength{\tabcolsep}{7pt}
\setlength{\emergencystretch}{2em}
\newcolumntype{Y}{>{\raggedright\arraybackslash}X}
\newcolumntype{R}{>{\raggedleft\arraybackslash}p{0.105\textwidth}}
\newcommand{\metric}{\textit{log loss}}
\newcommand{\code}[1]{\texttt{#1}}

\pagestyle{fancy}
\fancyhf{}
\setlength{\headheight}{14pt}
\lhead{PMLDL 2026 \textbar{} LLM Classification Fine-tuning}
\rhead{Technical Report}
\cfoot{\thepage}
\renewcommand{\headrulewidth}{0.4pt}

\begin{document}

\begin{center}
  {\LARGE\bfseries\color{accent} LLM Classification Fine-tuning}\\[0.35em]
  {\large Technical Report}\\[0.6em]
  {\small Анонимная версия для Stage 2}
\end{center}

\vspace{0.5em}

\section{Постановка задачи}

Для каждого примера известны пользовательский запрос и два ответа языковых
моделей. Необходимо предсказать предпочтение человека: ответ Модели A, ответ
Модели B или ничью. Использован набор данных соревнования Kaggle \cite{kaggle}.
Модель возвращает три вероятности в порядке: Модель A, Модель B, ничья.

Основной критерий качества --- multiclass \metric:

\begin{equation}
  \mathcal{L} = -\frac{1}{N}\sum_{i=1}^{N}\sum_{k=1}^{3}
  y_{ik}\log p_{ik},
  \label{eq:logloss}
\end{equation}

где $y_{ik}$ равно единице, если для $i$-го примера верен класс $k$, и нулю
иначе, а $p_{ik}$ --- предсказанная вероятность класса. Меньшее значение
означает лучшее качество и более точную калибровку вероятностей. Дополнительно
приведены accuracy, macro F1, ECE и Brier score; они не использовались для
выбора модели.

Все сравнения выполнялись по одному фиксированному протоколу, поэтому модели
не сравниваются на разных выборках. Folds 0--6 использовались только для
обучения, fold 7 --- для выбора кандидатов и абляций, fold 8 --- для подбора
весов ансамбля и калибровки. Финальный holdout fold 9 не использовался для
выбора гиперпараметров, состава ансамбля или температуры.

\begin{table}[htbp]
\centering
\caption{Фиксированное разбиение данных.}
\label{tab:protocol}
\small
\begin{tabularx}{\textwidth}{|>{\bfseries}p{0.26\textwidth}|p{0.17\textwidth}|R|Y|}
\hline
\rowcolor{lightblue} Часть данных & Folds & Примеров & Назначение \\ \hline
Обучение & 0--6 & 40\,235 & Обучение всех сравниваемых моделей. \\ \hline
Выбор кандидатов & 7 & 5\,746 & Baseline, HPO и выбор моделей для ансамбля. \\ \hline
Калибровка & 8 & 5\,748 & Подбор весов и temperature scaling после выбора кандидатов. \\ \hline
Финальная проверка & 9 & 5\,748 & Однократная независимая оценка замороженной конфигурации. \\ \hline
\end{tabularx}
\end{table}

\section{Baseline-модели}

Первая baseline-модель --- multinomial logistic regression на структурных
признаках диалога: длине prompt и ответов, разнице длин, числе ходов и простых
отношениях между ответами. Вторая модель использует TF-IDF признаков слов и
символьных $n$-грамм вместе с теми же структурными признаками. Для обеих
моделей применяются A/B swap augmentation при обучении и усреднение
предсказаний для исходного и переставленного порядка ответов.

Улучшенный sparse TF-IDF получен последовательным one-factor-at-a-time
поиском на fold 7. В финальной конфигурации используются word $n$-grams
1--3, char $n$-grams 2--5, словарь из 100\,000 word-признаков и параметр
$\alpha=3\cdot10^{-4}$ для линейного классификатора. Остальные части
протокола не менялись.

\begin{table}[htbp]
\centering
\caption{Baseline-результаты на общем selection fold 7.}
\label{tab:baselines}
\small
\begin{tabularx}{\textwidth}{|Y|R|R|R|R|}
\hline
\rowcolor{lightblue} Модель & \metric & Accuracy & Macro F1 & ECE \\ \hline
Structural Logistic Regression & 1.06303 & 0.44692 & 0.44014 & 0.02615 \\ \hline
Word + character TF-IDF & 1.04052 & 0.46920 & 0.46153 & 0.01621 \\ \hline
Улучшенный sparse TF-IDF & \textbf{1.03576} & \textbf{0.47355} & \textbf{0.46595} & 0.01879 \\ \hline
\end{tabularx}
\end{table}

Сочетание word и character TF-IDF снизило \metric{} на $0.02252$ относительно
структурной baseline. Последовательный sparse sweep дал дополнительное
снижение на $0.00476$. Sparse-модель уступает нейросетевой модели, однако её
ошибки достаточно отличаются от ошибок DeBERTa, поэтому она сохранена как
компонент ансамбля.

\section{Предлагаемая нейросетевая модель}

Основной нейросетевой кандидат --- \code{microsoft/deberta-v3-base}
\cite{deberta}, дообученный методом QLoRA \cite{qlora}. Вход строится из
prompt, ответа Модели A и ответа Модели B. Тексты сериализуются по ходам
диалога и ограничиваются balanced head-tail truncation: для длинного текста
сохраняются начало и конец каждого сегмента, а общий лимит равен 512 токенам.

При обучении ответы A и B меняются местами с вероятностью $0.5$; метка при
этом меняется согласованно. Во время inference вероятности для исходного и
переставленного примеров усредняются после обратной перестановки классов.
Это уменьшает зависимость результата от позиции ответа.

\begin{table}[htbp]
\centering
\caption{Конфигурация выбранной DeBERTa QLoRA модели.}
\label{tab:deberta-card}
\small
\begin{tabularx}{\textwidth}{|>{\bfseries}p{0.35\textwidth}|Y|}
\hline
\rowcolor{lightblue} Параметр & Значение \\ \hline
Backbone & \code{microsoft/deberta-v3-base}, revision \code{8ccc9b6f36199bec6961081d44eb72fb3f7353f3}. \\ \hline
Обучение & QLoRA в 4-bit NF4 с double quantization. \\ \hline
Адаптеры & $r=16$, $\alpha=32$, dropout $=0.05$; \code{query\_proj} и \code{value\_proj}. \\ \hline
Дополнительно обучаемые модули & \code{pooler} и \code{classifier}. \\ \hline
Оптимизатор & AdamW, learning rate $2.8\cdot10^{-4}$, weight decay $0.01$. \\ \hline
Обучение & 3 эпохи, effective batch size 32, warmup ratio 0.06. \\ \hline
Симметрия A/B & Swap augmentation и swap-time averaging. \\ \hline
\end{tabularx}
\end{table}

Выбранная модель с seed 42 получила на fold 7 \metric{} $1.01505$, accuracy
$0.48747$ и macro F1 $0.48420$. Она стала лучшим отдельным кандидатом среди
проверенных моделей.

\section{Улучшения и контролируемые эксперименты}

Для выполнения controlled improvement был проведён заранее зафиксированный
QLoRA hyperparameter search. Менялись только rank адаптера, dropout, learning
rate и выбираемая эпоха; backbone, разбиение, токенизация, оптимизатор и
A/B-протокол оставались одинаковыми. Все пять вариантов сравнивались по
primary metric на fold 7.

\begin{table}[htbp]
\centering
\caption{Подбор гиперпараметров DeBERTa QLoRA на fold 7.}
\label{tab:qlora-hpo}
\small
\begin{tabularx}{\textwidth}{|Y|R|R|R|R|R|}
\hline
\rowcolor{lightblue} Конфигурация & Rank & Dropout & Learning rate & Лучшая эпоха & \metric \\ \hline
Reference r16 & 16 & 0.05 & $2.0\cdot10^{-4}$ & 3 & 1.02507 \\ \hline
Low learning rate r16 & 16 & 0.05 & $1.2\cdot10^{-4}$ & 5 & 1.04791 \\ \hline
High learning rate r16 & 16 & 0.05 & $2.8\cdot10^{-4}$ & 3 & \textbf{1.01627} \\ \hline
Rank 32 & 32 & 0.05 & $2.0\cdot10^{-4}$ & 2 & 1.03805 \\ \hline
Rank 32, dropout 0.10 & 32 & 0.10 & $2.0\cdot10^{-4}$ & 1 & 1.08046 \\ \hline
\end{tabularx}
\end{table}

Лучший HPO-вариант снизил \metric{} на $0.00880$ относительно reference r16.
Для него fold-7 log loss после третьей эпохи составил $1.01627$, а после
четвёртой вырос до $1.01883$; поэтому применено раннее прекращение, а пятая
эпоха не использовалась. Повторный clean run выбранной конфигурации с seed 42
дал $1.01505$.

Также проверена чувствительность выбранной QLoRA-конфигурации к seed.
Усреднение трёх seed не улучшило primary metric, поэтому в ансамбль вошёл
только лучший одиночный seed 42.

\begin{table}[htbp]
\centering
\caption{Влияние seed на DeBERTa QLoRA, fold 7.}
\label{tab:seeds}
\small
\begin{tabularx}{\textwidth}{|Y|R|R|R|Y|}
\hline
\rowcolor{lightblue} Конфигурация & \metric & Accuracy & Macro F1 & Решение \\ \hline
Seed 17 & 1.02800 & 0.46711 & 0.46777 & Не выбран. \\ \hline
Seed 42 & \textbf{1.01505} & \textbf{0.48747} & \textbf{0.48420} & Выбран. \\ \hline
Seed 73 & 1.02807 & 0.47355 & 0.47328 & Не выбран. \\ \hline
Среднее трёх seed & 1.01889 & 0.48016 & 0.48001 & Уступает seed 42. \\ \hline
\end{tabularx}
\end{table}

\section{Результаты}

На fold 7 были выбраны DeBERTa QLoRA seed 42 и improved sparse TF-IDF.
На fold 8 сравнивались три фиксированных набора положительных весов. Для
каждого набора применялось multiclass temperature scaling. Лучшая комбинация
--- равное усреднение вероятностей двух моделей с температурой
$T=0.817311$. Public Kaggle score не использовался при выборе моделей,
весов или температуры.

\begin{table}[htbp]
\centering
\caption{Сравнение отдельных моделей и ансамбля на calibration fold 8.}
\label{tab:calibration}
\small
\begin{tabularx}{\textwidth}{|Y|R|R|R|R|R|}
\hline
\rowcolor{lightblue} Модель & \metric & Accuracy & Macro F1 & ECE & Brier score \\ \hline
Калиброванная DeBERTa QLoRA & 1.00892 & 0.48713 & 0.48349 & 0.01833 & 0.60422 \\ \hline
Калиброванный sparse TF-IDF & 1.02609 & 0.49791 & 0.49214 & 0.03329 & 0.61043 \\ \hline
Финальный ансамбль & \textbf{0.99697} & \textbf{0.50035} & \textbf{0.49521} & \textbf{0.01111} & \textbf{0.59635} \\ \hline
\end{tabularx}
\end{table}

\begin{table}[htbp]
\centering
\caption{Перебор весов после temperature scaling на fold 8.}
\label{tab:weights}
\small
\begin{tabularx}{\textwidth}{|Y|>{\raggedleft\arraybackslash}p{0.15\textwidth}|>{\raggedleft\arraybackslash}p{0.15\textwidth}|>{\raggedleft\arraybackslash}p{0.15\textwidth}|}
\hline
\rowcolor{lightblue} Веса (DeBERTa, sparse) & Температура & \metric & ECE \\ \hline
0.25, 0.75 & 0.87399 & 1.00397 & 0.01542 \\ \hline
0.50, 0.50 & \textbf{0.81731} & \textbf{0.99697} & \textbf{0.01111} \\ \hline
0.75, 0.25 & 0.84529 & 0.99899 & 0.01746 \\ \hline
\end{tabularx}
\end{table}

Равный ансамбль улучшил \metric{} на $0.01195$ относительно калиброванной
DeBERTa и на $0.02911$ относительно калиброванной sparse-модели. Снижение
ECE до $0.01111$ показывает, что итоговые вероятности лучше согласованы с
наблюдаемыми частотами классов. Так как fold 8 использовался для калибровки,
его результат рассматривается как подтверждение выбранной конфигурации, а
не как независимая оценка обобщающей способности.

\section{Воспроизводимость}

Исходный код, фиксированные folds, конфигурации, результаты и инструкции
находятся в проекте. Пакеты зафиксированы в
\code{requirements-baseline.lock} и \code{requirements-transformer.lock}.
Финальная конфигурация ансамбля с весами и температурой хранится в
\code{configs/final\_model.json}.

\begin{table}[htbp]
\centering
\caption{Артефакты, необходимые для повторения итогового pipeline.}
\label{tab:artifacts}
\small
\begin{tabularx}{\textwidth}{|>{\bfseries}p{0.36\textwidth}|Y|}
\hline
\rowcolor{lightblue} Артефакт & Назначение \\ \hline
\code{data/splits/folds.csv} и \code{metadata.json} & Неизменное распределение примеров по folds. \\ \hline
\code{data/checksums.sha256} & Проверка версии официальных Kaggle CSV. \\ \hline
\code{artifacts/models/}\\\code{deberta-v3-base/} & Offline snapshot backbone точной revision \code{8ccc9b6f...}. \\ \hline
\code{artifacts/models/}\\\code{qlora\_adapter.zip} & Обученный QLoRA adapter seed 42; SHA-256 указан в manifest. \\ \hline
\code{artifacts/models/}\\\code{sparse\_estimator.joblib} & Выбранный sparse TF-IDF estimator; SHA-256 \code{5be06b0efdd...}. \\ \hline
\code{configs/final\_model.json} & Классовый порядок, веса $0.5/0.5$ и температура $0.817311$. \\ \hline
\code{docs/REPRODUCE.md} & Порядок установки окружения, подготовки данных и запуска. \\ \hline
\end{tabularx}
\end{table}

Минимальный запуск начинается с создания виртуального окружения, установки
двух lock-файлов и команды \code{make audit}. Она проверяет контрольные суммы
официальных данных и соответствие зафиксированным folds. Затем отдельные
scripts воспроизводят sparse sweep, QLoRA search, калибровку и финальное
усреднение без изменения параметров из \code{final\_model.json}. Полные
команды приведены в \code{docs/REPRODUCE.md}.

\section{Заключение}

В проекте сравнены несколько baseline-моделей, дообучена DeBERTa-v3-base
методом QLoRA и проведён контролируемый поиск её гиперпараметров. Лучший
отдельный кандидат --- DeBERTa QLoRA seed 42 с fold-7 log loss $1.01505$.
Улучшенный sparse TF-IDF слабее как самостоятельная модель, но дополняет
нейросетевой кандидат в ансамбле.

Выбранное решение --- равное усреднение вероятностей DeBERTa QLoRA и sparse
TF-IDF с temperature scaling $T=0.817311$. На calibration fold оно достигло
log loss $0.99697$, accuracy $0.50035$, macro F1 $0.49521$ и ECE $0.01111$.
Главные ограничения --- выбор на одном selection fold, вариативность между
seed и необходимость сохранять неизменной конфигурацию до независимой
проверки на fold 9.

\begin{thebibliography}{9}

\bibitem{kaggle}
Kaggle. \emph{LLM Classification Fine-tuning}.\\
\url{https://www.kaggle.com/competitions/llm-classification-finetuning}

\bibitem{deberta}
P. He, X. Liu, J. Gao, and W. Chen. \emph{DeBERTa: Decoding-enhanced BERT
with Disentangled Attention}. ICLR, 2021.\\
\url{https://arxiv.org/abs/2006.03654}

\bibitem{lora}
E. Hu et al. \emph{LoRA: Low-Rank Adaptation of Large Language Models}.
ICLR, 2022.\\
\url{https://arxiv.org/abs/2106.09685}

\bibitem{qlora}
T. Dettmers et al. \emph{QLoRA: Efficient Finetuning of Quantized LLMs}.
NeurIPS, 2023.\\
\url{https://arxiv.org/abs/2305.14314}

\end{thebibliography}

\end{document}
